\documentclass[aps,prd,twocolumn,10pt,nofootinbib,superscriptaddress]{revtex4-2}
\usepackage{amsmath,amssymb,graphicx}
\usepackage{booktabs}
\usepackage[colorlinks=true,linkcolor=blue,citecolor=blue,urlcolor=blue]{hyperref}

\newcommand{\aeff}{a_0}
\newcommand{\FF}{\mathcal{F}}
\newcommand{\Msun}{M_\odot}

\begin{document}

\title{A relativistic completion for the dark-energy--anchored acceleration scale
$a_0=\kappa c\sqrt{G\rho_\Lambda}$:\\
aether--scalar--tensor theory with an offset-DBI khronon and an $a_0$-resonant environment response}

\author{Carl P.\ Zimmerman}
\affiliation{Briar Creek Tech}
\date{2026-08-09}

\begin{abstract}
We assemble and stress-test a relativistic field theory carrying the acceleration scale
$a_0=\kappa c\sqrt{G\rho_\Lambda}=9.36\times10^{-11}\,\mathrm{m\,s^{-2}}$, where $\rho_\Lambda$ is the
measured dark-energy density and $\kappa$ is a pure number that we treat as \emph{measured}, not
derived: a distance-free estimator built on the SPARC sample gives $\kappa=0.551\pm0.043$
($\pm0.063$ if the $H_0$ tension is carried as a systematic), consistent with $\tfrac12$ but not
singling it out. The host is aether--scalar--tensor (AeST) theory [Skordis \& Z\l o\'snik, PRL
\textbf{127}, 161302 (2021)], the only relativistic MOND-class theory known to reproduce the CMB;
the content of this work is the free function, which factorises into three terms with distinct jobs.
(i) A MOND kernel normalised by $a_0$, convex and bijective, with deep-MOND and Newtonian limits
exact and a Newtonian residual of order $e^{-\sqrt{g/a_0}}$ ($\sim10^{-3457}$ at Earth's orbit).
Because AeST has no gravitational slip ($\Phi=\Psi$, $\gamma_{\rm PPN}=1$), dynamical and lensing
masses agree identically, converting the $21.2\sigma$ lensing exclusion that kills modified-inertia
realisations into a $0.6\sigma$ pass. (ii) An offset Dirac--Born--Infeld (DBI) function of the
khronon's temporal invariant, whose minimum supplies $w=-1$ exactly (dark energy), whose quadratic
deviations supply a pressureless cosmological component, and whose boundedness we show is
\emph{necessary}: for any power-law $K\sim u^n$ the early-time equation of state limits to
$w=1/(n-1)$, so only a bounded $K$ reaches dust. This reverses the sign of the published $455\times$
tension between cosmological and galactic requirements on the Helmholtz mass [Blanchet \& Skordis,
JCAP 11 (2024) 040], turning cosmology into a lower bound aligned with MOND ($226\times$ margin).
A Boltzmann-code (CLASS) confrontation leaves TT unchanged to $0.069\%$ and $P(k=0.2\,h/\mathrm{Mpc})$
to $1.7\%$. (iii) An environment response $A\,B(Y/a_0^2)(Q-Q_0)^2$ with $B(y)=y/(1+y)^2$: a
position-dependent Helmholtz mass peaked at the theory's own $a_0$, motivated by the observation that
cluster interiors are the unique deep-potential environments sitting at the MOND transition
($R_{500}$ at $0.33$--$0.58\,a_0$). The term passes a five-environment matrix (clusters by
calibration; galaxy interiors at $3.4\times10^{-4}$ of the baryonic mass; galaxy outskirts at 1 Mpc
at $1.2\%$ against a $5.9\%$ weak-lensing bound; the linear cosmos at $0.6\%$ of mean matter; the
solar system at $10^{-19}$), vanishes identically on FRW, and enters linear cosmology at second
order. A full perturbation analysis including all aether mixings yields: a no-ghost \emph{theorem}
(the term's kinetic contribution is rank-one positive semi-definite, so by Weyl monotonicity it
cannot lower any eigenvalue of the healthy AeST kinetic matrix); $c_T=1$ exactly; a derived
quasi-static closure ($\delta A^0=-\Phi A^0$); and two-sided amplitude bounds that leave the
mechanism falsifiable from both directions. The theory uses five dark-sector parameters against
$\Lambda$CDM's two; what is bought is that dark energy, the cosmological dark-matter phenomenology,
and MOND arise from one function. We list eight explicit non-claims, led by ``$\kappa=\tfrac12$ is
not derived'' and including the sharpest open problem: whether the dust component stays out of
galaxy interiors, as the data demand, awaits a nonlinear relaxation calculation never performed for
this theory class. Six falsifiable predictions follow, led by a hash-stamped, pre-registered Gaia
DR4 wide-binary target. Every quantitative claim is reproducible from named, committed scripts that exit
non-zero on failure.
\end{abstract}

\maketitle

\section{Introduction}
\label{sec:intro}

Two facts anchor this paper. First, galactic dynamics exhibits a characteristic acceleration: below
$g\sim10^{-10}\,\mathrm{m\,s^{-2}}$ rotation curves flatten, and the observed centripetal
acceleration follows the baryonic prediction through a universal interpolation --- the radial
acceleration relation (RAR) \cite{McGaugh2016,Lelli2016} --- with remarkably small scatter. Milgrom
identified the regularity in 1983 and codified it as MOND \cite{Milgrom1983}. Second, the measured
value of that scale is numerically close to accelerations constructible from the cosmological
constant, as Milgrom also observed. The present program takes the second fact literally and asks
what follows if
\begin{equation}
a_0 \;=\; \kappa\, c\sqrt{G\rho_\Lambda} \;=\; \frac{cH_\Lambda}{Z} \;=\;
c^2\sqrt{\frac{\Lambda}{32\pi}} \;=\; 9.3619\times10^{-11}\,\mathrm{m\,s^{-2}},
\label{eq:a0}
\end{equation}
with $\kappa=\tfrac12$, $Z=2\sqrt{8\pi/3}$, and $H_\Lambda$ the asymptotic de~Sitter expansion rate,
is a statement about gravity rather than a coincidence. The three forms in Eq.~(\ref{eq:a0}) are
algebraically identical; $\kappa$ and $Z$ are conventions on one number.

A one-line relation is not a theory. It cannot bend light, form structure, or be checked for
internal consistency. This paper assembles the relativistic field theory that carries
Eq.~(\ref{eq:a0}) --- ``the completion'' --- and reports the outcome of an extended internal
adversarial program: uniqueness and no-go theorems, environment matrices, stability analyses, a
Boltzmann-code confrontation, and the public withdrawal of three failed cluster mechanisms that
preceded the surviving one. Throughout, results that could have killed the theory are given the same
prominence as results that support it, and every quantitative claim is backed by a named script in
the public repository \cite{repo} that recomputes it from scratch and exits non-zero on failure
(Appendix~\ref{app:repro}).

Three disclosures frame everything that follows. (1) The host theory, aether--scalar--tensor (AeST),
is Skordis and Z\l o\'snik's \cite{SZ2021}; our content is the free function and the theorems about
it. (2) An earlier arm of this program realised Eq.~(\ref{eq:a0}) through modified inertia; it is
\emph{excluded} at $21.2\sigma$ by the equality of dynamical and lensing masses and was withdrawn.
The present modified-gravity realisation is the surviving arm. (3) $\kappa$ is measured, not
derived; Sec.~\ref{sec:coefficient} quantifies exactly how well, and Sec.~\ref{sec:notclaimed}
lists what is not claimed.

\section{The scale: form, interpolation, coefficient}
\label{sec:coefficient}

\subsection{The form is forced}

If $a_0$ is built from $(G,c,\rho)$ alone, dimensional analysis admits exactly one combination,
$a_0=\xi c\sqrt{G\rho}$: the exponent matrix is nonsingular (determinant 2), so the form is unique,
and the theorem has content because admitting $\hbar$ destroys the uniqueness. The choice
$\rho=\rho_\Lambda$ is the physical hypothesis; the form then follows.

\subsection{The interpolation is derivable --- and its coefficient is excluded}

A de~Sitter--Unruh construction sharpens the stakes. Using the Deser--Levin temperature of an
accelerated observer in de~Sitter space \cite{DeserLevin} and Milgrom's 1999 vacuum-heuristic
\cite{Milgrom1999}, the MOND interpolation is \emph{fully derived}, with no fitted element:
$\nu(y)=\sqrt{1+1/y}$, which is precisely Milgrom's Eq.~(9) of Ref.~\cite{Milgrom1999}. The same
construction, however, is rigid at every step --- the power is forced by the two MOND limits, the
baseline by the ambient Gibbons--Hawking temperature, the normalisation by the Newtonian limit ---
and it forces $a_0=2cH_\Lambda$, which the SPARC data exclude at $15.6\sigma$. The construction that
derives the \emph{function} cannot supply the \emph{coefficient}; the heuristic must never be cited
as support for $\kappa=\tfrac12$. Hence the coefficient must be measured.

\subsection{The coefficient, measured}

Two routes. The baryonic Tully--Fisher intercept on SPARC gives $\hat\kappa=0.465\pm0.076$, with a
hard $9.47\%$ mass-budget floor (stellar and gas calibrations trade against each other at fixed
$f_\star+f_{\rm gas}=1$). A sharper, distance-free estimator exploits the fact that
$g_{\rm bar}\propto F/\theta^2$ (flux over angular size squared) carries no distance dependence
while $g_{\rm obs}\propto1/D$ does; built on the $g_{\rm bar}$ axis of the RAR it gives
\begin{equation}
\kappa \;=\; 0.551\pm0.043 \quad (\pm0.063~\text{under the } H_0~\text{tension}),
\end{equation}
since $\kappa\propto a_0/H_0$ at full strength. Four simple candidates
($\tfrac12$, $1/\sqrt3$, $\sqrt{3/8}$, $0.40$) sit inside $2\sigma$: the data are consistent with
$\tfrac12$ and do not single it out. The binding systematics are the absolute 21-cm flux scale and
CO-dark molecular gas, not stellar populations; an irreducible $\sim3.9\%$ floor on $a_0$ caps the
$\tfrac12$-vs-$1/\sqrt3$ discrimination near $4\sigma$ for any SPARC-class analysis.

\section{The action}
\label{sec:action}

We adopt the AeST action of Ref.~\cite{SZ2021}: gravity is Einstein--Hilbert; matter couples to the
metric only; the dark sector is a unit-timelike vector $A^\mu$ ($A^\mu A_\mu=-1$) and a scalar
$\varphi$ entering through the two invariants
\begin{equation}
Q \equiv A^\mu\nabla_\mu\varphi, \qquad
Y \equiv (g^{\mu\nu}+A^\mu A^\nu)\nabla_\mu\varphi\,\nabla_\nu\varphi .
\end{equation}
The content of this work is the free function, which factorises:
\begin{equation}
\boxed{\;\FF(Y,Q) \;=\; \frac{a_0^2}{8\pi G}\,\FF_Y\!\Big(\frac{Y}{a_0^2}\Big)
\;+\; K(Q) \;+\; A\,B\!\Big(\frac{Y}{a_0^2}\Big)(Q-Q_0)^2\;}
\label{eq:F}
\end{equation}
with $B(y)=y/(1+y)^2$. The three terms are treated in Secs.~\ref{sec:galaxies},
\ref{sec:cosmology}, and \ref{sec:clusters} respectively.

\emph{The Y-sector.} $\FF_Y$ is fixed by the kernel $\nu(y)=1/(1-e^{-\sqrt y})$,
$y=g_{\rm bar}/a_0$ --- the functional form of the empirical RAR fitting function of
Ref.~\cite{McGaugh2016} --- equivalently the closed parametric pair
\begin{equation}
\mu(u)=1-e^{-u},\qquad x(u)=\frac{u^2}{\mu(u)},\qquad u=\sqrt y,\; x=\frac{g_{\rm obs}}{a_0},
\label{eq:kernel}
\end{equation}
normalised by Eq.~(\ref{eq:a0}).

\emph{The Q-sector.} An offset-DBI function,
\begin{equation}
K(Q) = -M^4 + \mu^2\Lambda_D^2\Big[1-\sqrt{1-\tfrac{u^2}{\Lambda_D^2}}\Big],
\quad u\equiv Q-Q_0,\;|u|<\Lambda_D,
\label{eq:K}
\end{equation}
with $M^4=\rho_\Lambda$ the vacuum scale, $\mu$ the Helmholtz mass, $\Lambda_D$ the field-space
cap, and the conserved shift charge $I_0$ setting the dust abundance.

\emph{The environment response.} The third term is a position-dependent Helmholtz mass
$\mu^2_{\rm eff}=A\,B(g^2/a_0^2)$ peaked at the theory's own $a_0$; the peak location costs nothing
($Y/a_0^2$ is already the Y-sector's normalisation), and the single new calibrated amplitude is
$A\approx1.7\,\mathrm{Mpc}^{-2}$.

\section{Galactic phenomenology}
\label{sec:galaxies}

The kernel~(\ref{eq:kernel}) embeds exactly: the deep-MOND limit $\mu\to x$ and the Newtonian limit
$\mu\to1$ are exact, $x(u)$ is a bijection ($h'=e^{-u}(1+u)>0$), and the resulting free function is
convex, from which the baryonic Tully--Fisher relation follows exactly. The Newtonian residual is
$e^{-\sqrt y}$: at Earth's orbit, $3.6\times10^{-3457}$ --- planetary ephemerides are untouched at
any conceivable precision. On the RAR, the anchored $a_0$ of Eq.~(\ref{eq:a0}) fits SPARC at
$0.108$ dex scatter at $\Upsilon_\star=0.70$, inside the population-synthesis range (computed with
the framework's Milgrom-1999-shape interpolation, which is data-degenerate with the kernel of
Eq.~(\ref{eq:kernel}): their divergence is below the scatter everywhere). Stated honestly: the same
estimator at the literature's fitted $a_0=1.2\times10^{-10}$ with $\Upsilon_\star$ frozen at $0.5$
gives $0.122$ dex, but with $\Upsilon_\star$ refit on both sides the two are statistically
indistinguishable --- $a_0$ and $\Upsilon_\star$ are degenerate on the RAR, so the relation does
not by itself discriminate the conventions. What anchoring buys is not a better fit but a cheaper
one: the scale is taken from an external cosmological measurement rather than from the galaxies
being explained, at no cost in fit quality.

The environment response also operates inside galaxies and must not spoil this: it contributes
$3.4\times10^{-4}$ of the baryonic mass in galaxy interiors, shifting the RAR by
$4\times10^{-4}$ dex --- two orders below the observed scatter (Sec.~\ref{sec:clusters}). A separate
internal requirement --- that the Q-sector dust of Sec.~\ref{sec:cosmology} not accumulate in galaxy
interiors --- is open, and is stated as such in Sec.~\ref{sec:notclaimed}, non-claim 2.

\section{Lensing}
\label{sec:lensing}

AeST has no gravitational slip: the weak-field potentials obey $\Phi=\Psi$, so
$\gamma_{\rm PPN}=1$ and $M_{\rm dyn}/M_{\rm lens}=1$ identically, for any source. This is the
structural fact that separates the modified-gravity realisation of Eq.~(\ref{eq:a0}) from the
modified-inertia realisation: the latter predicts $M_{\rm dyn}/M_{\rm lens}=1/f_{\rm bar}\simeq6.4$
and is excluded at $21.2\sigma$; the former passes at $0.601\sigma$. The dark field genuinely
sources curvature where it accumulates, so light cannot distinguish this theory from a dark-matter
universe at the level of the deflection law.

\section{Cosmology: the offset-DBI sector}
\label{sec:cosmology}

\subsection{Three regimes of one function}

At the condensate minimum $u=0$: $\rho=-K=M^4$, $p=K=-M^4$, hence $w=-1$ \emph{exactly} --- dark
energy is the resting value of the khronon's energy function, proved rather than assumed. For
$0<|u|\ll\Lambda_D$: $K\to\mu^2u^2/2$ and the excitation dilutes as $a^{-3}$ with vanishing
pressure --- the cosmological dust. As $|u|\to\Lambda_D$: $K$ is bounded and $K'\to\infty$, and the
early-time equation of state returns to $w\to0$.

\subsection{Boundedness is necessary: a no-go for power laws}

For any power-law $K\sim u^n$ the early-time equation of state limits to
\begin{equation}
w \;\longrightarrow\; \frac{1}{\,n-1\,},
\end{equation}
so the quadratic gives a stiff $w\to1$ phase, the quartic $w\to\tfrac13$, the sextic $w\to\tfrac15$:
every polynomial fails big-bang nucleosynthesis and CMB constraints in the same calculable way.
Only a bounded $K$ reaches $w\to0$. This theorem re-frames the published $455\times$ conflict
between the cosmological and galactic requirements on $\mu^{-1}$ \cite{BlanchetSkordis2024}, derived
there for the quadratic form: with the DBI form the early stiff phase is absent and the sign of the
cosmological requirement \emph{reverses} --- cosmology becomes a lower bound on $\mu^{-1}$, aligned
with the MOND-side requirement, with a $226\times$ margin. The no-go is not evaded but dissolved:
its premise (the quadratic) is the one choice the theorem forbids. The DBI sector is ghost-free
($K''=\mu^2(1-s^2)^{-3/2}>0$ on the whole range) and subluminal
($c_s^2\le0.385\,\Lambda_D$). The cost is disclosed: an EFT boundary at saturation
($K''/\mu^2\sim3\times10^{11}$ at the cap), where the description is effective only.

\subsection{Boltzmann-code confrontation}

A real CLASS run with the DBI sector mapped to a generalised-dark-matter fluid gives, at
$\Lambda_D\le10^{-2}$ (natural units of Ref.~\cite{SZ2021}): $c_s^2(z_{\rm rec})=2.9\times10^{-8}$,
TT spectrum unchanged to $0.069\%$, and $P(k=0.2\,h/\mathrm{Mpc})$ to $1.71\%$ --- indistinguishable
from CDM on the primary CMB. Two caveats are part of the result. First, the pass depends on the
sound-speed bump being transient (it peaks at $z\approx189$); a patched fluid carrying the full
$c_s^2(a)$ history through structure growth is required and owed. Second, the stability analysis of
Sec.~\ref{sec:health} subsequently forced $\Lambda_D\le8.4\times10^{-7}$, \emph{excluding the
fiducial value used in the CLASS run}; since the run's agreement only improves with smaller
$\Lambda_D$ (the sound speed falls), the pass survives, but the window
$1.9\times10^{-10}\ll\Lambda_D\le8.4\times10^{-7}$ (3.7 decades) is now health-bounded from above.

\subsection{What the CMB requires, exactly}

The CMB requires a component that gravitates, is pressureless, and dilutes as $a^{-3}$, at the full
$\Omega_{\rm dm}$: deleting it moves the third-to-first acoustic peak ratio by $54\%$ and no refit
absorbs it ($\Delta\chi^2>400$). It does not require a particle: by the generalised-dark-matter
degeneracy \cite{Kopp2018}, a fluid excitation and a particle species are indistinguishable at
$0\sigma$ in the linear CMB. This theory supplies the component as the khronon's shift-charge
excitation. The defensible statement --- and the only one this paper makes --- is: \emph{no
dark-matter particle; the cosmological pressureless job is done by the same field that supplies
$\Lambda$ and MOND.} The flat claim ``there is no dark matter'' is not available to this theory or
to any theory that matches the CMB. Nor, at present, is ``none in galaxies'': the galaxy data are
fit by baryons plus the kernel alone, so the theory \emph{needs} its dust component to stay out of
galaxy interiors, and whether it does is an open nonlinear problem (Sec.~\ref{sec:notclaimed},
non-claim 2).

\section{Clusters: the $a_0$-resonant environment response}
\label{sec:clusters}

\subsection{Why the obvious routes fail}

MOND under-predicts cluster masses; any completion must supply the residual without disturbing
galaxies that share the relevant scales. Three mechanisms were built and killed within this
program, and the kills shape the survivor:

\emph{(a) The $R^2$ lever} (residual density tracking the potential, ratio fixed by system size)
requires $\mu^{-1}=3.13$ Mpc, violating both galaxy-outskirt weak-lensing bounds of
Ref.~\cite{Mistele2023}. Withdrawn.

\emph{(b) The primordial shift-charge route} (clusters khronon-rich because their Lagrangian
patches were over-dense) fails three independent ways, the third decisive: (i) cluster $R_{500}$
and galaxy-outskirt lensing sample wavenumbers only $1.4\times$ apart ($k\simeq4.5$ vs
$6.3\,\mathrm{Mpc}^{-1}$) with disjoint abundance requirements; (ii) for correlated Gaussian
initial conditions $\mathbb{E}[\delta_{\rm kh}|\delta_m]$ is linear in $\delta_m$, so no such
statistic can select environments; (iii) \emph{smooth accretion}: any suppression mechanism removes
fluctuations, not the mean, and a collapsing region captures all cold content in its basin
(the khronon obeys the identical shell equation by the equivalence principle), so
$\xi(\mathrm{halo})\to1$ for \emph{any} cold transfer function. Withdrawn, and the prior version of
this paper's cluster section with it.

\emph{(c) The potential response} $\rho_c\propto\Phi^n$: cosmological potentials are scale-free
($\Phi_{\rm rms}\sim10^{-5}$ everywhere, only $2.2\times$ below cluster centres), so any $n$ steep
enough to select clusters floods the linear cosmos at $\sim20\times$ the mean matter density.
Closed.

The lesson: neither scale nor depth alone selects clusters. What does is the pair (depth,
$g\sim a_0$): cluster interiors at $R_{500}$ sit at $0.33$--$0.58\,a_0$ --- \emph{clusters are the
cosmic objects that live at the MOND transition}.

\subsection{The term, and the environment matrix}

The third term of Eq.~(\ref{eq:F}) implements exactly that selection: a Helmholtz mass
$\mu^2_{\rm eff}=A\,B(g^2/a_0^2)$ with $B(y)=y/(1+y)^2$ peaked at $g=a_0$, falling on both sides.
Calibrating $A\approx1.7\,\mathrm{Mpc}^{-2}$ on the cluster residual
($\mu^2_{\rm eff}=0.23\,\mathrm{Mpc}^{-2}$), the term then passes the four non-calibration
environments: galaxy interiors, $3.4\times10^{-4}$ of $M_b$ ($B$ near peak but the potential
shallow); galaxy outskirts at 1 Mpc, $1.2\%$ of the baryonic mass against the $5.9\%$ strict bound
from Ref.~\cite{Mistele2023} stacking; the linear cosmos, $0.6\%$ of mean matter; the solar system,
$10^{-19}$. On FRW, $Y=0$ and the term vanishes identically; it enters perturbations at second
order, so the linear CMB and $P(k)$ of Sec.~\ref{sec:cosmology} are untouched by construction. It
evades kill (b) because a response has no charge to advect, and kill (c) because it does not couple
to depth alone. Against interest: cluster modelling in Ref.~\cite{Mistele2023}'s own framework
prefers amplitudes $4$--$34\times$ our chain calibration; the health bounds below cut into exactly
that band, making the tension a two-sided test rather than a tuning freedom.

\subsection{Stability, and the full perturbation matrix}
\label{sec:health}

\emph{No ghosts, as a theorem.} The term's contribution to the kinetic matrix is
$2AB(y)\ge0$ in the $\chi\chi$ entry alone --- a rank-one positive semi-definite addition. By Weyl's
monotonicity theorem, adding a PSD matrix cannot lower any eigenvalue; AeST's kinetic matrix is
healthy in its stable region; therefore the completed theory is ghost-free wherever AeST is,
\emph{for all field values and all environments}. Verified in closed form on the $2\times2$
subsystem and numerically on random healthy $4\times4$ matrices. There is no Ostrogradsky sector
(no higher derivatives are introduced).

\emph{Tensor speed.} The new terms are algebraic in the metric perturbation: $c_T=1$ exactly, and
GW170817 \cite{GW170817} is satisfied by construction, not by tuning.

\emph{Quasi-static closure.} The unit-norm constraint gives $\delta A^0=-\Phi A^0$, hence
$\delta Q=\dot\chi-Q_0\Phi$; the full matrix regenerates the quasi-static phenomenology assumed in
the environment matrix. The gradient block including all scalar--aether mixings is, exactly,
\begin{equation}
\Delta G \;=\; 2\lambda
\begin{pmatrix} q & 2y^{3/2}B'' \\ 2y^{3/2}B'' & \,y\,q \end{pmatrix},
\qquad
q=\frac{3y^2-8y+1}{(1+y)^4},
\end{equation}
with $\lambda$ the dimensionless bump amplitude in the Y-sector's gradient normalisation (exact
assembly in the scripts). $q<0$ precisely on $0.13<y<2.5$ --- the cluster band --- which is what
turns stability into a cap on $A$.

\emph{The bounds.} On FRW, gradient health forces $\Lambda_D\le8.4\times10^{-7}$, overruling this
program's own earlier fiducial $10^{-2}$; the $\Phi$-mixing correction is Poisson-suppressed
($3\times10^{-4}$) and does not reopen it. In halos, the isolated-term analysis caps
$A\le2.7\,A_{\rm fid}$; integrating out the aether softens this to
$A_{\rm max}=(2.7\text{--}7.4)\,A_{\rm fid}$, the residual range pinned by one coefficient of the
quasi-static expansion of Ref.~\cite{SZ2021}'s supplemental material (a literature lookup). Either
way, the demanding ($34\times$) end of the Ref.~\cite{Mistele2023} band is excluded by
$>7\times$ and the $4\times$ edge is marginal: the amplitude is pinched between an observational
floor and a stability ceiling. The fiducial calibration sits inside but near the edge
($1+2\lambda q=+0.63$ at the cluster minimum). Two flags remain open and are so labelled: a
$\lambda$-suppressed tachyon-type vector mass in $y>1$ interiors (watch item), and the known
$k<\mu$ formal unboundedness of AeST \cite{SZ2021}, which at this theory's $\mu^{-1}=4392$ Mpc
migrates to the Hubble scale (within $15\%$) --- relocated, not resolved.

\emph{Free signature.} The response carries anisotropic stress of order $0.6$ in clusters: a
specific lensing-vs-hydrostatic mass discrepancy, testable on existing samples
(Sec.~\ref{sec:predictions}).

\emph{Status.} The mechanism is a \emph{candidate}: it has survived this program's environment
matrix, health check, and full perturbation analysis, and no external scrutiny. A real cluster-model
fit is owed.

\section{Parameter count, honestly}
\label{sec:parameters}

$\Lambda$CDM's dark sector: $\Omega_c,\Omega_\Lambda$. This completion: $M$ ($=\rho_\Lambda^{1/4}$),
$I_0$, $\mu$, $\Lambda_D$, $A$ --- \emph{five numbers against two}, with $a_0$ not independent:
$a_0=m_{\rm cond}/(4\sqrt\pi)$, $m_{\rm cond}=M^2/(\sqrt2 M_{\rm Pl})$, holds to $0.076\%$ --- but
see non-claim (1) below for why this is a relabelling rather than a derivation. What the extra
parameters buy is structural: $\Lambda$, the cosmological dust, and MOND from one function instead
of three unrelated sectors, plus kill-tests $\Lambda$CDM does not offer.

\section{What is not claimed}
\label{sec:notclaimed}

\begin{enumerate}
\item \textbf{$\kappa=\tfrac12$ is not derived.} $a_0=m_{\rm cond}/(4\sqrt\pi)$ is algebraically
identical to $a_0=\tfrac12 c\sqrt{G\rho_\Lambda}$ and to $Z=2\sqrt{8\pi/3}$; the residue is
convention-dependent (reduced vs non-reduced Planck mass), and a derived number cannot depend on
bookkeeping. The best current framing is a graviton-bath calculation in which the de~Sitter horizon
entropy cancels the Planck suppression exactly ($S_{\rm dS}\,GH^2=\pi$ identically), leaving
$\kappa^2=8\pi\,\epsilon_{\rm tot}$ with $\epsilon_{\rm tot}$ a pure number from horizon-mode
counting --- well-posed and open. Five defensible normalisations span $\kappa=0.013$--$2.047$; the
one landing exactly on $\tfrac12$ is structurally invalid ($h_{\mu\nu}u^\mu u^\nu$ vanishes in TT
gauge for a static worldline) and is not cited.
\item \textbf{``No dark matter'' is not claimed --- and ``none in galaxies'' is an open problem,
not a result.} The pressureless component exists at full $\Omega_{\rm dm}$
(Sec.~\ref{sec:cosmology}D); what does not exist is a particle. Whether the dust stays out of galaxy
interiors is genuinely open: the smooth-accretion theorem that killed cluster route (b) implies it
falls into every collapsing basin, galaxies included. The favorable branch --- relaxation to the
Helmholtz-preferred, centrally-evacuated static profile ($\rho_c$ tracking the potential,
$\xi\propto R^2$, galaxy-interior cost $\sim6\times10^{-5}$ dex) --- and the fatal branch ---
CDM-like concentration, overshooting the RAR by the banked $2.06$--$4.42\times$ --- are both
computable; the deciding quantity is the condensate's relaxation time toward the static profile
within a Hubble time, a nonlinear calculation never performed for AeST by any group. It is this
program's next scheduled computation. Until then, galaxy-interior cleanliness is an assumption the
theory needs, not a result it owns.
\item \textbf{The cluster mechanism is a candidate}, pinched and falsifiable, not a result.
\item \textbf{The post-recombination growth history is partially verified.} The acoustic-peak pass
is real; a patched $c_s^2(a)$ fluid through structure growth is owed, as is a Lyman-$\alpha$
confrontation of the dust component's small-scale initial conditions.
\item \textbf{Cassini $Q_2$ is inherited} ($3$--$15\sigma$ class tension) and not relieved here.
\item \textbf{PPN $\alpha_1,\alpha_2$} for this free function are not yet computed against
lunar-laser-ranging and binary-pulsar bounds.
\item \textbf{The wide-binary target is the point-field asymptote}; the full nonlinear AQUAL-EFE
solve is owed before Gaia DR4.
\item \textbf{This is not a theory of everything}; nothing here touches the Standard Model. The
2026-06-23 retraction of earlier overclaims stands.
\end{enumerate}

\section{Falsifiable predictions}
\label{sec:predictions}

\begin{enumerate}
\item \emph{Wide binaries (Gaia DR4).} Registered, hash-stamped target
$\gamma_v=\sqrt{\nu(y_{\rm ext,N})}=1.2139$ (canonical footing; $1.2592$ alternative), against the
withdrawn modified-inertia arm's $1.1582$ --- separated at $2.68\sigma$ by the frozen pipeline's own
error budget. The pre-registration and nine hash-stamped amendments are committed \cite{repo};
DR4-as-frozen measures $a_0$ to $12.6\%$.
\item \emph{Directional external-field asymmetry.} AQUAL-class prediction $1$--$4\%$, signed; pure
modified inertia predicted exactly zero. A first firing on existing data gives $\hat A=+2.95$,
$p=0.029$, with the predicted sign --- suggestive, not decisive.
\item \emph{Cluster anisotropic stress.} Order-$0.6$ lensing-vs-hydrostatic mass discrepancy with a
specific profile --- the bump's free signature.
\item \emph{Cluster-to-cluster scatter} in the residual, $\sigma(\xi)/\xi\approx0.02$--$0.60$,
distinguishing the environment response from mechanisms predicting zero scatter, on existing data.
\item \emph{$a_0(z)$.} If $a_0\propto\sqrt{\rho_{\rm de}(z)}$, then
$a_0(z)/a_0(0)=(1+z)^{\frac32(1+w_0+w_a)}e^{-\frac32 w_a z/(1+z)}$ --- a bump-then-decline under
DESI-class $(w_0,w_a)$, distinct from constancy and from $cH(z)$ scaling.
\item \emph{No particle, ever.} Any confirmed direct detection of a dark-matter particle falsifies
the theory outright.
\end{enumerate}

\section{Conclusion}

The relation $a_0=\kappa c\sqrt{G\rho_\Lambda}$ now has a relativistic carrier: AeST with the
factorised free function of Eq.~(\ref{eq:F}). Within it, lensing, the primary CMB, the RAR, the
BTFR, and the solar system are quantitatively passed; dark energy is the condensate minimum rather
than an input; the published quadratic-sector no-go is dissolved by a boundedness theorem; and the
long-standing cluster problem has, for the first time in this program, a candidate mechanism that
survived its own execution squad --- resonant, falsifiably, at the theory's own acceleration scale.
The theory is more expensive than $\Lambda$CDM in parameters and more constrained in structure; its
distinctive wager is that one measured coincidence is a law. The wager is now fully specified,
stability-certified, and scheduled for execution or vindication by data that exist or are imminent.

\appendix

\section{Reproducibility}
\label{app:repro}

Every quantitative claim above is recomputed by a committed script in Ref.~\cite{repo}
(\texttt{real\_research/reviews/}), each of which exits non-zero on failure and carries negative
controls. Key mappings: lensing and the AeST embedding,
\texttt{mi\_relativistic\_completion\_aest\_2026.py}; the DBI sector, its theorem and health,
\texttt{mi\_dbi\_khronon\_2026.py}; the CLASS run, \texttt{mi\_dbi\_cmb\_class\_run\_2026.py}; the
CMB dust requirement, \texttt{mi\_cmb\_no\_dust\_existence\_2026.py}; the condensate--$a_0$
relation, \texttt{mi\_condensate\_vacuum\_energy\_a0\_2026.py}; the Deser--Levin derivation and its
$15.6\sigma$ scissors, \texttt{mi\_deser\_levin\_interpolation\_2026.py}; the coefficient
measurement, \texttt{mi\_distance\_free\_gbar\_estimator\_sparc\_2026.py} and
\texttt{mi\_kappa\_error\_budget\_unlock\_2026.py}; the cluster-route kills,
\texttt{mi\_ic\_route\_1mpc\_confrontation\_2026.py}; the environment response, its health, and the
full perturbation matrix, \texttt{mi\_a0\_bump\_response\_2026.py},
\texttt{mi\_a0\_bump\_health\_2026.py}, \texttt{mi\_aest\_full\_matrix\_bump\_2026.py}; the
graviton-bath framing of $\kappa$, \texttt{mi\_graviton\_bath\_ctp\_2026.py}; kernel embedding and
BTFR, \texttt{mi\_route\_a\_field\_theory\_2026.py}.

\begin{thebibliography}{99}

\bibitem{Milgrom1983} M.~Milgrom, Astrophys.\ J.\ \textbf{270}, 365 (1983).

\bibitem{Milgrom1999} M.~Milgrom, Phys.\ Lett.\ A \textbf{253}, 273 (1999).

\bibitem{BekensteinMilgrom1984} J.~Bekenstein and M.~Milgrom, Astrophys.\ J.\ \textbf{286}, 7 (1984).

\bibitem{Bekenstein2004} J.~D.~Bekenstein, Phys.\ Rev.\ D \textbf{70}, 083509 (2004).

\bibitem{SZ2021} C.~Skordis and T.~Z\l o\'snik, Phys.\ Rev.\ Lett.\ \textbf{127}, 161302 (2021).

\bibitem{McGaugh2016} S.~S.~McGaugh, F.~Lelli, and J.~M.~Schombert, Phys.\ Rev.\ Lett.\
\textbf{117}, 201101 (2016).

\bibitem{Lelli2016} F.~Lelli, S.~S.~McGaugh, and J.~M.~Schombert, Astron.\ J.\ \textbf{152}, 157
(2016).

\bibitem{DeserLevin} S.~Deser and O.~Levin, Class.\ Quantum Grav.\ \textbf{14}, L163 (1997).

\bibitem{ArkaniHamed2004} N.~Arkani-Hamed, H.-C.~Cheng, M.~A.~Luty, and S.~Mukohyama, JHEP
\textbf{0405}, 074 (2004).

\bibitem{BlanchetSkordis2024} L.~Blanchet and C.~Skordis, JCAP \textbf{11} (2024) 040.

\bibitem{Verwayen2024} P.~Verwayen, C.~Skordis, and C.~Boehm, Mon.\ Not.\ R.\ Astron.\ Soc.\
\textbf{531}, 272 (2024).

\bibitem{Durakovic2024} A.~Durakovic and C.~Skordis, JCAP \textbf{04} (2024) 040.

\bibitem{Mistele2023} T.~Mistele, S.~McGaugh, and S.~Hossenfelder, Astron.\ Astrophys.\
\textbf{676}, A100 (2023).

\bibitem{Kopp2018} M.~Kopp, C.~Skordis, D.~B.~Thomas, and S.~Ili\'c, Phys.\ Rev.\ Lett.\
\textbf{120}, 221102 (2018).

\bibitem{GW170817} B.~P.~Abbott \emph{et al.}, Astrophys.\ J.\ Lett.\ \textbf{848}, L13 (2017).

\bibitem{RogersPeiris2021} K.~K.~Rogers and H.~V.~Peiris, Phys.\ Rev.\ Lett.\ \textbf{126}, 071302
(2021).

\bibitem{McGaugh2024} S.~S.~McGaugh, J.~M.~Schombert, F.~Lelli, and J.~Franck, Astrophys.\ J.\
(2024), accelerated structure formation and JWST high-$z$ massive galaxies.

\bibitem{repo} C.~P.~Zimmerman, the public research repository,
\url{https://github.com/carlzimmerman/zimmerman-formula}; technical companion: \emph{The
Completion}, DOI \href{https://doi.org/10.5281/zenodo.21863521}{10.5281/zenodo.21863521};
plain-language companion, same repository.

\end{thebibliography}

\end{document}
